\section{Introduction}
\label{sec:intro}

Enterprise collections contain regulations, failure analyses and operating
procedures. When an agent starts from a catalogue, a vague description can
send it to the wrong collection. Even when the search scope is already
known, a stale index can point to a document's former path. Both failures
come from the state the retrieval system reads, not from the passage ranker.

Document acquisition, organisation and retrieval are often studied
separately. Enterprise content management studies document
governance \cite{paivarinta-ecm-2005,tyrvaainen-ecm-2006}, while metadata
quality and tagging affect how users discover documents
\cite{park-metadata-quality-2009,golder-tagging-2006}. Distributed
information retrieval addresses the selection of collections
\cite{callan1995searching,si2003redde}. Retrieval-augmented generation
(RAG) uses ranked passages to inform generation, and modular retrieval
workflows can expose evidence through interfaces such as the Model
Context Protocol (MCP) \cite{rag-lewis,dpr,modular-rag,mcp,gao-rag-survey-2023}.
This paper asks how a retrieval workflow can use managed document state
and how the platform exposes that state for operator inspection when
relevant evidence is missed.

The Hierarchical Management Retrieval System (HMRS) is a self-hosted
platform that links document management to evidence access. During
ingestion, the workflow records where a document belongs and what it
contains. Search uses this state to select a collection and document or
to start within a specified collection. The system exposes collection
listings, document descriptions and source references so that the stages
that supplied the candidates can be inspected.

\noindent\textbf{Contributions.}
\begin{enumerate}
\item \textbf{A management model for evidence access.} HMRS connects
three managed levels (collection, document and segment) to four access
views (catalogue, shelf, document and segment). Collection membership
sets the search scope, descriptions guide selection and source references
link passages to documents (Section~\ref{sec:system}).
\item \textbf{A content-informed document lifecycle with post-operation
probes.} The acquisition workflow derives placement, descriptions and
tags from parsed content. Operators can reconcile files, tree entries
and library metadata after lifecycle changes and probe retrieval for
stale index entries
(Sections~\ref{sec:acquisition}--\ref{sec:organization}).
\item \textbf{Agent access policies over shared state.} MCP tools let an
agent navigate from collection and document descriptions to source text,
or search the indices of specified collections directly
(Section~\ref{sec:retrieval}).
\item \textbf{An evidence gate that keeps admission and ordering as
separate, inspectable decisions.} An optional server-side gate fetches
candidate text, scores passages and returns those that meet the threshold
with source references, reporting fetch and scoring failures separately;
the caller keeps control of candidate order and can distinguish an
incomplete review from a completed review that finds no support
(Section~\ref{sec:gate}). The separation is load-bearing rather than
stylistic: ordering candidates by the admission scores alone displaced
gold documents from the top-five window in the replay, so the two
decisions are exposed, measured and reported separately
(Section~\ref{sec:e6}).
\end{enumerate}

The evaluation checks three links in this design. First, we ask whether
changes to a managed document are visible to search. Next, we test whether
the selector reaches the labelled collection, whether the evidence pack
contains the required text, and whether the verifier's score is suitable
for ordering candidates. Finally, we assess answer quality in the native configurations. The pinned
HMRS configuration in the answer comparison did not exercise catalogue
navigation, and its gate calls did not succeed (Section~\ref{sec:e10}).
The component experiments address those functions separately.

\subsection{Research questions and reading guide}
\label{sec:rq}

\looseness=-1 \textbf{RQ1} examines document acquisition, collection membership and
lifecycle consistency. \textbf{RQ2} measures collection selection under
catalogue-first access. \textbf{RQ3} compares target-string containment
and retrieval text volume across the tested evidence configurations.
\textbf{RQ4} asks how verifier-score ordering changes gold-document
Hit@5 in fixed candidate pools. \textbf{RQ5} examines abstention and
operational failures. \textbf{RQ6} compares judged answer quality under
native configurations.

Section~\ref{sec:related} positions HMRS within related work.
Sections~\ref{sec:system}--\ref{sec:implementation} describe the management
model and its implementation. Section~\ref{sec:eval} reports the evaluation,
followed by the discussion, limitations and conclusions in
Sections~\ref{sec:discussion}--\ref{sec:conclusion}.
